%% kapa: exam=1 round=37 year=2026 period=2 form=A subject=accounting subject_ko=회계학 seq=38 number=78
%% source: Q-Net 감정평가사 기출문제 (cst00309) artlSeq=5258924 file=raw/1008/5258924/2026년 제37회 감정평가사 1차 2교시.pdf
%% conversion: kordoc 4.24.51 -> tools/gen1.py (auto); review=proofread-sonnet; flags=diag-header-approx
\begin{question}{78}
(주)감평은 두 개의 제조부문(P1, P2)과 두 개의 보조부문(S1, S2)을 운영하고 있으며, 상호배분법을 사용하여 보조부문원가를 제조부문에 배분한다. 각 부문 간의 용역수수관계는 다음과 같다.
\par
{\small\centering
\begin{tabular}{|c|c|c|c|c|}\hline
\multirow{2}{*}{\makecell{\hfill 사용부문\\제공부문\hfill}} & \multicolumn{2}{c|}{보조부문} & \multicolumn{2}{c|}{제조부문} \\\cline{2-5}
 & S1 & S2 & P1 & P2 \\\hline
S1 & - & 30 \% & 50 \% & 20 \% \\\hline
S2 & 15 \% & 50 \% & 20 \% & 15 \% \\\hline
\end{tabular}\par}
두 개의 보조부문(S1, S2)으로부터 제조부문 P1과 P2에 배분된 금액이 각각 ￦142,500과 ￦78,000인 경우, 보조부문원가를 배분하기 이전의 보조부문 S1과 S2에 집계된 원가는? \par
\hspace*{2.4em}\makebox[4.5cm][l]{\uline{S1}}\uline{S2} \par
\begin{choices}
  \item \makebox[4.5cm][l]{￦100,000,}￦120,500
  \item \makebox[4.5cm][l]{￦100,500,}￦120,000
  \item \makebox[4.5cm][l]{￦120,000,}￦100,500
  \item \makebox[4.5cm][l]{￦150,000,}￦165,000
  \item \makebox[4.5cm][l]{￦165,000,}￦150,000
\end{choices}
\end{question}
\answer{3}
